\documentclass[12pt,a4paper]{ctexart} % 使用 ctexart 支持中英双语

\usepackage{amsmath, amssymb, bm}
\usepackage{geometry}
\usepackage{indentfirst}
\usepackage{hyperref}
\usepackage{graphicx}

\geometry{margin=1in}
\setlength{\parindent}{2em}

\title{%
\textbf{误差状态卡尔曼滤波器（ESKF）设计文档}\\[0.2cm]
\large Error-State Kalman Filter (ESKF) Design for Underwater Navigation
}
\author{Underwater Robot Navigation Project}
\date{}

\begin{document}
\maketitle

\begin{abstract}
\textbf{中文：}  
本文档详细阐述了本项目中用于水下机器人导航定位的误差状态卡尔曼滤波（ESKF）建模与推导方法，包括名义状态、误差状态、系统动力学、离散化预测模型、DVL 测量模型、噪声特性以及统一时间基设计。本设计同时作为 C++ 实时导航模块与 Python 离线分析的共同数学基础。

\textbf{English:}  
This document presents the complete mathematical formulation of the 
Error-State Kalman Filter (ESKF) used in the underwater robot navigation 
system. It includes the nominal and error-state definitions, 
continuous-time dynamics, discrete-time propagation, DVL measurement 
model, noise specifications, and unified timebase convention. 
This design serves as the mathematical foundation for both the C++ 
real-time navigation module and Python offline analysis.
\end{abstract}

\section{Overview / 概述}

\textbf{中文：}  
ESKF 将系统状态拆分为“名义状态”和“误差状态”，在预测阶段使用非线性动力学推进名义状态，在更新阶段基于线性化误差状态进行卡尔曼校正，从而获得高精度与高稳定性。此方法特别适合 IMU+低频传感器（如 DVL）的组合导航。

\textbf{English:}  
The ESKF splits the system into a nonlinear nominal state propagated by 
IMU integration and a linear error-state used for Kalman updates. 
This structure provides numerical stability and accuracy, 
especially suitable for IMU + low-rate sensors such as DVL.

\section{State Definitions / 状态定义}

\subsection{Nominal State (16D) / 名义状态（16维）}

\[
\mathbf{x} =
\begin{bmatrix}
\mathbf{p} \\
\mathbf{v} \\
\mathbf{q} \\
\mathbf{b_a} \\
\mathbf{b_g}
\end{bmatrix}
\]

中文说明：  
\begin{itemize}
  \item \(\mathbf{p}\)：位置（ENU）
  \item \(\mathbf{v}\)：速度（ENU）
  \item \(\mathbf{q}\)：姿态四元数（ENU$\rightarrow$Body）
  \item \(\mathbf{b_a}, \mathbf{b_g}\)：IMU 加计与陀螺偏置
\end{itemize}

English explanation:  
\begin{itemize}
  \item \(\mathbf{p}\)：position in ENU
  \item \(\mathbf{v}\)：velocity
  \item \(\mathbf{q}\)：attitude quaternion
  \item \(\mathbf{b_a}, \mathbf{b_g}\)：accelerometer and gyro biases
\end{itemize}

\subsection{Error State (15D) / 误差状态（15维）}

\[
\delta \mathbf{x} =
\begin{bmatrix}
\delta\mathbf{p} \\
\delta\mathbf{v} \\
\delta \boldsymbol{\theta} \\
\delta\mathbf{b_a} \\
\delta\mathbf{b_g}
\end{bmatrix}
\]

中文：姿态误差用小角度 \(\delta\boldsymbol{\theta}\) 表示。\\
English: Attitude error is represented using a small-angle vector \(\delta\boldsymbol{\theta}\).

\[
\delta\mathbf{q} \approx 
\begin{bmatrix}1 \\ \frac{1}{2}\delta\boldsymbol{\theta}\end{bmatrix}
\]

\section{IMU Continuous Dynamics / IMU 连续时间动力学}

IMU 测量模型：
\[
\mathbf{a}_m = \mathbf{a} + \mathbf{b_a} + \mathbf{n_a},\qquad
\boldsymbol{\omega}_m = \boldsymbol{\omega} + \mathbf{b_g} + \mathbf{n_g}.
\]

\subsection{Nominal Position / 名义位置}
\[
\dot{\mathbf{p}} = \mathbf{v}.
\]

\subsection{Nominal Velocity / 名义速度}
\[
\dot{\mathbf{v}} = \mathbf{R}(\mathbf{q})(\mathbf{a}_m - \mathbf{b_a}) + \mathbf{g}.
\]

\subsection{Nominal Attitude / 名义姿态}
\[
\dot{\mathbf{q}} = \frac{1}{2} \Omega(\boldsymbol{\omega}_m - \mathbf{b_g}) \mathbf{q}.
\]

\subsection{Bias Drift / 偏置随机游走}
\[
\dot{\mathbf{b_a}} = \mathbf{n_{ba}},
\qquad
\dot{\mathbf{b_g}} = \mathbf{n_{bg}}.
\]

\section{Error-State Dynamics / 误差状态线性化模型}

\subsection{Position Error / 位置误差}
\[
\delta\dot{\mathbf{p}} = \delta\mathbf{v}.
\]

\subsection{Velocity Error / 速度误差}
\[
\delta \dot{\mathbf{v}} =
-\mathbf{R}[\mathbf{a}_m - \mathbf{b_a}]_\times \delta\boldsymbol{\theta}
-\mathbf{R}\delta\mathbf{b_a}
+ \mathbf{R}\mathbf{n_a}.
\]

\subsection{Attitude Error / 姿态误差}
\[
\delta \dot{\boldsymbol{\theta}} =
-(\boldsymbol{\omega}_m - \mathbf{b_g})_\times \delta\boldsymbol{\theta}
- \delta \mathbf{b_g}
+ \mathbf{n_g}.
\]

\subsection{Bias Errors / 偏置误差}
\[
\delta \dot{\mathbf{b_a}} = \mathbf{n_{ba}},
\qquad
\delta \dot{\mathbf{b_g}} = \mathbf{n_{bg}}.
\]

\section{Discrete-Time Prediction / 离散预测模型}

IMU 频率高（50--200 Hz），使用 Euler 离散化足够：

\subsection{Nominal propagation / 名义状态推进}
\[
\mathbf{p}_{k+1} = \mathbf{p}_{k} + \mathbf{v}_{k} dt,
\]
\[
\mathbf{v}_{k+1} = \mathbf{v}_{k}
+ \mathbf{R}(\mathbf{a}_m - \mathbf{b_a}) dt
+ \mathbf{g} dt,
\]
\[
\mathbf{q}_{k+1} = \mathbf{q}_k \otimes \delta\mathbf{q},
\qquad
\delta\mathbf{q} =
\begin{bmatrix}
1 \\ \frac{1}{2}(\boldsymbol{\omega}_m - \mathbf{b_g})dt
\end{bmatrix}.
\]

\subsection{Covariance propagation / 协方差推进}
\[
\mathbf{P}_{k+1} =
\mathbf{F} \mathbf{P}_k \mathbf{F}^T
+ \mathbf{G}\mathbf{Q}\mathbf{G}^T.
\]

\section{DVL Measurement Model / DVL 测量模型}

DVL 直接测量 ENU 速度：
\[
\mathbf{z}_{dvl} =
\begin{bmatrix}
v_E \\ v_N \\ v_U
\end{bmatrix}
= \mathbf{v} + \mathbf{n_v}.
\]

\subsection{Measurement Function / 测量函数}
\[
\mathbf{h}(\mathbf{x}) = \mathbf{v}.
\]

\subsection{Jacobian / 观测矩阵}
\[
\mathbf{H} =
\begin{bmatrix}
\mathbf{0}_{3\times3} &
\mathbf{I}_{3\times3} &
\mathbf{0}_{3\times9}
\end{bmatrix}.
\]

\subsection{Kalman Update / 卡尔曼更新}

创新：
\[
\mathbf{y} = \mathbf{z}_{dvl} - \hat{\mathbf{v}}.
\]

创新协方差：
\[
\mathbf{S} = \mathbf{H} \mathbf{P} \mathbf{H}^T + \mathbf{R}.
\]

增益：
\[
\mathbf{K} = \mathbf{P} \mathbf{H}^T \mathbf{S}^{-1}.
\]

误差：
\[
\delta\mathbf{x} = \mathbf{K}\mathbf{y}.
\]

\section{Error Injection / 误差注入}

位置和速度：
\[
\mathbf{p} \leftarrow \mathbf{p} + \delta\mathbf{p},
\qquad
\mathbf{v} \leftarrow \mathbf{v} + \delta\mathbf{v}.
\]

姿态：
\[
\delta\mathbf{q} =
\begin{bmatrix}
1 \\ \frac{1}{2}\delta\boldsymbol{\theta}
\end{bmatrix},
\qquad
\mathbf{q} \leftarrow \delta\mathbf{q} \otimes \mathbf{q}.
\]

偏置：
\[
\mathbf{b_a} \leftarrow \mathbf{b_a} + \delta\mathbf{b_a},
\qquad
\mathbf{b_g} \leftarrow \mathbf{b_g} + \delta\mathbf{b_g}.
\]

协方差：
\[
\mathbf{P} \leftarrow (\mathbf{I} - \mathbf{K}\mathbf{H})\mathbf{P}.
\]

\section{Noise and Covariance Models / 噪声模型与协方差选择}

IMU 噪声：
\[
\mathbf{Q} = \mathrm{diag}(
\sigma_a^2, \sigma_g^2, \sigma_{ba}^2, \sigma_{bg}^2
).
\]

DVL 噪声：
\[
\mathbf{R}_{dvl} = \mathrm{diag}(\sigma_v^2).
\]

说明：如果 DVL bottom lock 质量不好，应增大 $\sigma_v$ 并进行门控。

\section{Timebase / 时间基统一机制}

系统使用：
\begin{itemize}
    \item \textbf{MonoNS：}单调时钟，供 ESKF 内部计算 $dt$
    \item \textbf{EstNS：}系统时钟，用于日志与人类可读时间
\end{itemize}

时间差：
\[
dt = \frac{\Delta \mathrm{MonoNS}}{10^9}.
\]

\section{Conclusion / 总结}

本文档提供了：
\begin{itemize}
  \item 完整的 ESKF 数学推导与符号
  \item 统一的噪声模型与时间基定义
  \item C++ 实时导航模块与 Python 离线分析的共同数学依据
\end{itemize}

This document serves as the unified mathematical reference for all 
contributors working on ESKF, real-time navigation, data processing, 
and underwater robot dynamics modeling.

\end{document}
